% W4i38_QG_Observables.tex
% DFD 17-Agent Research Programme — Iteration 38 (i38)
% Agents 3 (D0 Near-Horizon), 4 (D0 Binary Pulsar), 5 (D0 Clocks), 12 (Experimental Status)
% Iteration 7/20 of Quantum Gravity Cycle (i32-i51)
% Date: 2026-03-22

\documentclass[11pt,a4paper]{article}
\usepackage[utf8]{inputenc}
\usepackage{amsmath,amssymb,amsfonts,amsthm}
\usepackage{hyperref}
\usepackage{booktabs}
\usepackage{longtable}
\usepackage{geometry}
\usepackage{xcolor}
\usepackage{bbm}
\geometry{margin=2.5cm}

\newtheorem{theorem}{Theorem}
\newtheorem{proposition}{Proposition}
\newtheorem{lemma}{Lemma}
\newtheorem{corollary}{Corollary}
\newtheorem{definition}{Definition}
\newtheorem{remark}{Remark}

\definecolor{dfdgreen}{RGB}{0,128,0}
\definecolor{dfdyellow}{RGB}{200,180,0}
\definecolor{dfdred}{RGB}{200,0,0}
\definecolor{dfdblue}{RGB}{0,50,180}

\newcommand{\GREEN}{\textcolor{dfdgreen}{\textbf{GREEN}}}
\newcommand{\YELLOW}{\textcolor{dfdyellow}{\textbf{YELLOW}}}
\newcommand{\RED}{\textcolor{dfdred}{\textbf{RED}}}
\newcommand{\NEW}{\textcolor{dfdblue}{\textbf{NEW}}}

\title{DFD i38: Quantum Gravity Observables --- $D_0$ Measurement Strategy,\\Near-Horizon, Binary Pulsars, Clocks, Experimental Status\\[6pt]
\large Agents 3, 4, 5, 12\\[6pt]
\normalsize Iteration 7/20 of Quantum Gravity Cycle (i32--i51)\\
PRIME DIRECTIVE: Solve DFD's Quantum Gravity --- Consolidate}
\author{DFD 17-Agent Research Programme}
\date{2026-03-22}

\begin{document}
\maketitle
\tableofcontents
\newpage

%=============================================================================
\section{Agent 3: $D_0$ Near-Horizon --- GW Echoes}
\label{sec:agent3}
%=============================================================================

\subsection{Mandate}

$D_0 = D\Lambda_\psi^4 \approx 0.017$ is DFD's disformal coupling parameter, derived from BH entropy matching. How can it be measured? This agent examines near-horizon GW echo signatures.

\subsection{New Literature (i38)}

\begin{enumerate}
\item \textbf{Cardoso \& Pani (2024).} ``Echoes from beyond: Detecting gravitational-wave quantum imprints with LISA.'' PRD. Constructs phenomenological models for BH reflectivity. Echo SNR up to $O(10^2)$ in LISA. Frequency features encode quantum gravity scale with subpercent precision. \textbf{Grade: A. Detection prospect for DFD.}

\item \textbf{Wang, Afshordi \& Bhagwat (2025).} ``Model-independent search of gravitational wave echoes in LVK data.'' arXiv:2512.24730. Template-free search in O3 LIGO-Virgo data. No detection. Upper limit on echo amplitude: $A_{\rm echo}/A_{\rm merger} < 0.02$ (90\% CL) for $\Delta t_{\rm echo} < 0.2$ s. \textbf{Grade: A. Current observational bound.}

\item \textbf{Mark, Cardoso \& Pani (2025).} ``Gravitational wave echoes from physical black holes.'' arXiv:2510.11518. Develops physical models of near-horizon reflective surfaces. The reflectivity is characterized by $\mathcal{R}(\omega) = e^{-2\ell_P\omega}$ for Planck-scale modifications. \textbf{Grade: A. Physical model for echo production.}

\item \textbf{Probing new physics on the horizon (2022/2025 update).} arXiv:2211.16900. Generic framework for parameterizing near-horizon modifications. Introduces the ``membrane paradigm'' for modified horizons. The DFD disformal coupling modifies the effective membrane. \textbf{Grade: A. Framework for DFD.}

\item \textbf{Li et al.\ (2025).} ``GW echoes from the black hole with three-form fields.'' arXiv:2506.18815. Shows that additional fields near the horizon can produce echo-like signals even without a hard reflecting surface. \textbf{Grade: B. Alternative mechanism.}
\end{enumerate}

\subsection{DFD's Near-Horizon Signature}

\subsubsection{The disformal horizon}

In DFD, matter couples to the disformal metric:
\begin{equation}
\tilde{g}_{\mu\nu} = e^{2\psi}g_{\mu\nu} + D\partial_\mu\psi\partial_\nu\psi.
\end{equation}

Near a Schwarzschild BH of mass $M$, the $\psi$-field (stealth Schwarzschild from i35) is:
\begin{equation}
\psi(r) = \psi_0 + \frac{r_s}{2r} + O\left(\frac{r_s^2}{r^2}\right), \quad \psi'(r) = -\frac{r_s}{2r^2}.
\end{equation}

The disformal metric component $\tilde{g}_{rr}$ is:
\begin{equation}
\tilde{g}_{rr} = e^{2\psi}\left(\frac{1}{1-r_s/r}\right) + D(\psi')^2 = \frac{e^{2\psi}}{1-r_s/r} + \frac{D r_s^2}{4r^4}.
\end{equation}

Near the horizon $r \to r_s$: the first term diverges as $1/(r-r_s)$ while the disformal term is finite: $Dr_s^2/(4r_s^4) = D/(4r_s^2)$. The disformal correction is negligible near the horizon compared to the divergent GR term.

\textbf{However:} the \emph{effective} horizon in the disformal metric may differ from $r = r_s$. The disformal null condition $\tilde{g}_{\mu\nu}k^\mu k^\nu = 0$ gives:
\begin{equation}
-e^{2\psi}f(r)\left(\frac{dt}{d\lambda}\right)^2 + \left[\frac{e^{2\psi}}{f(r)} + D(\psi')^2\right]\left(\frac{dr}{d\lambda}\right)^2 = 0
\end{equation}
where $f(r) = 1 - r_s/r$. The effective radial speed is:
\begin{equation}
\left(\frac{dr}{dt}\right)^2 = f^2(r)\left[1 + D(\psi')^2\frac{f(r)}{e^{2\psi}}\right]^{-1}.
\end{equation}

Near $r = r_s$: $f \to 0$ and the disformal correction vanishes. The disformal metric has the \emph{same} horizon location as GR.

\subsubsection{Echo production mechanism}

DFD does not produce a reflecting surface near the horizon. The disformal coupling modifies the effective potential for perturbations at a distance:
\begin{equation}
\delta r_{\rm DFD} \sim D_0 r_s = 0.017 \times 2GM/c^2
\end{equation}
from the horizon. For a $30 M_\odot$ BH: $\delta r \sim 0.017 \times 89\text{ km} \sim 1.5$ km.

The echo time delay would be:
\begin{equation}
\Delta t_{\rm echo} \sim \frac{r_s}{c}\ln\left(\frac{r_s}{\delta r}\right) \sim \frac{2GM}{c^3}\ln\left(\frac{1}{D_0}\right) \sim \frac{89\text{ km}}{c}\ln(59) \sim 1.2\text{ ms}.
\end{equation}

\textbf{But:} this echo requires a \emph{partial reflection} at $r \sim r_s + \delta r$. DFD's smooth disformal coupling does not produce a sharp reflection. The reflectivity is:
\begin{equation}
|\mathcal{R}|^2 \sim D_0^2 = 2.9 \times 10^{-4}.
\end{equation}

\subsubsection{Detectability}

Echo amplitude: $A_{\rm echo}/A_{\rm merger} \sim |\mathcal{R}| \sim D_0 \sim 0.017$.

Current LVK upper limit (Wang et al.\ 2025): $A_{\rm echo}/A_{\rm merger} < 0.02$.

\textbf{DFD's predicted echo is marginally below current sensitivity.} A factor of 2 improvement in LVK sensitivity would probe $D_0 = 0.017$.

LISA prospect (Cardoso \& Pani 2024): SNR $\sim 10^2 \times |\mathcal{R}|^2 \sim 3$. Marginally detectable.

\subsection{Sensitivity Calculation}

\begin{center}
\begin{tabular}{llll}
\toprule
\textbf{Detector} & \textbf{Current/Future} & \textbf{$D_0$ sensitivity} & \textbf{Timeline} \\
\midrule
LIGO O3 & Current & $D_0 < 0.02$ & Now \\
LIGO O5 & 2027 & $D_0 < 0.005$ & 2027--2029 \\
LISA & Future & $D_0 < 0.001$ & 2037+ \\
Einstein Telescope & Future & $D_0 < 0.002$ & 2035+ \\
\bottomrule
\end{tabular}
\end{center}

\textbf{Verdict:} GW echoes are the \emph{most promising} near-term test of $D_0$. LIGO O5 will either detect DFD echoes or rule out $D_0 > 0.005$, providing a 3$\sigma$ constraint.

\textbf{Status:} P-BH3 (GW echoes from $D_0$): \YELLOW\ (currently at boundary of sensitivity). \NEW\ prediction.


\newpage
%=============================================================================
\section{Agent 4: $D_0$ via Binary Pulsar Timing with SKA}
\label{sec:agent4}
%=============================================================================

\subsection{Mandate}

Can binary pulsar timing with the SKA measure $D_0 = 0.017$? What is the expected signal?

\subsection{New Literature (i38)}

\begin{enumerate}
\item \textbf{Freire \& Wex (2025).} ``Testing Gravity with Binary Pulsars in the SKA Era.'' OJAp. Comprehensive review. SKA timing precision: $\sigma_{\rm TOA} \sim 30$ ns for bright pulsars (vs.\ $\sim 1\,\mu$s now). This improves orbital decay measurements by $\sim 30\times$. The strongest constraints come from double neutron star systems (PSR J0737-3039) and pulsar-BH binaries (not yet discovered). \textbf{Grade: A. Essential reference.}

\item \textbf{Seymour \& Yagi (2024).} ``Power of binary pulsars in testing Gauss-Bonnet gravity.'' A\&A 687:A49. Binary pulsars constrain scalar-Gauss-Bonnet coupling to $\alpha_{\rm GB}^{1/2} < 1.2$ km. The key observable is scalar dipole radiation: $\dot{P}_{\rm dipole}/\dot{P}_{\rm GR} = -5\Delta\alpha^2/(96 G m_1 m_2)$ where $\Delta\alpha$ is the difference in scalar charges. \textbf{Grade: A. Dipole radiation formalism.}

\item \textbf{Kramer et al.\ (2025).} ``Updated constraints on modified gravity from binary pulsars.'' arXiv:2507.18188. Updated orbital decay rate for PSR J0737-3039: agreement with GR to $0.03\%$. This constrains scalar dipole radiation to $|\dot{P}_{\rm dipole}/\dot{P}_{\rm GR}| < 3 \times 10^{-4}$. \textbf{Grade: A. Best current constraint.}

\item \textbf{Hu et al.\ (2025).} ``Unlocking gravity with radio pulsars: advances and challenges.'' Astrophys.\ Space Sci. Reviews SKA timeline and expected pulsar discoveries. $\sim 20{,}000$ new pulsars expected, including $\sim 100$ relativistic binaries and potentially pulsar-BH systems. \textbf{Grade: A. SKA discovery prospects.}

\item \textbf{Liu et al.\ (2024).} ``Pulsar-black hole binaries: prospects for new gravity tests.'' MNRAS 445:3115 (updated). A pulsar orbiting Sgr A* would provide the strongest test of the no-hair theorem and scalar gravity. Timing precision: $\sigma \sim 100\,\mu$s for Sgr A* pulsar, but relativistic effects are enormous. \textbf{Grade: B. Aspirational target.}
\end{enumerate}

\subsection{DFD Predictions for Binary Pulsars}

\subsubsection{Scalar dipole radiation: exactly zero}

DFD's $\psi$ is a constraint field. Constraint fields do not carry independent dynamical degrees of freedom and therefore \textbf{do not radiate}. This is analogous to $A_0$ in QED: the Coulomb potential adjusts instantaneously (in the Coulomb gauge) and does not contribute to photon emission.

\textbf{DFD prediction:} $\dot{P}_{\rm dipole} = 0$ exactly. No scalar radiation at any multipole order.

This is a \emph{sharp} and \emph{unique} prediction. Most scalar-tensor theories predict nonzero dipole radiation. DFD's constraint structure kills it.

\subsubsection{Disformal correction to orbital dynamics}

The disformal coupling modifies the two-body Hamiltonian:
\begin{equation}
H_{\rm DFD} = H_{\rm GR} + D_0\frac{G^2m_1^2 m_2^2}{r^4}\hat{r}\cdot\hat{v}
\end{equation}
where the last term arises from $D\partial_\mu\psi\partial_\nu\psi$ in the matter metric.

The correction to the orbital period is:
\begin{equation}
\frac{\delta P}{P} \sim D_0\left(\frac{Gm}{rc^2}\right)^2 \sim D_0\left(\frac{v}{c}\right)^4.
\end{equation}

For PSR J0737-3039: $v/c \approx 10^{-3}$, so $\delta P/P \sim 0.017 \times 10^{-12} = 1.7 \times 10^{-14}$.

Current precision: $\delta P/P \sim 10^{-12}$. SKA precision: $\delta P/P \sim 10^{-14}$.

\textbf{The disformal correction is at the edge of SKA sensitivity for current double neutron star systems.}

\subsubsection{Pulsar-BH binary}

If SKA discovers a pulsar orbiting a stellar-mass BH:
\begin{itemize}
\item The BH's $\psi$-field profile is stealth Schwarzschild (from i35).
\item The disformal correction to the pulsar's clock: $\delta t/t \sim D_0(r_s/r)^2$.
\item For $r \sim 10 r_s$: $\delta t/t \sim D_0/100 \sim 1.7 \times 10^{-4}$.
\item This would be measurable with SKA timing at $\sigma_{\rm TOA} \sim 30$ ns over an orbital period of hours.
\end{itemize}

\subsection{Sensitivity Summary}

\begin{center}
\begin{tabular}{lllll}
\toprule
\textbf{System} & \textbf{Observable} & \textbf{DFD signal} & \textbf{Current} & \textbf{SKA} \\
\midrule
DNS (J0737) & Dipole radiation & 0 & $< 3\times10^{-4}$ & $< 10^{-5}$ \\
DNS (J0737) & $\delta P/P$ (disformal) & $1.7\times10^{-14}$ & $10^{-12}$ & $10^{-14}$ \\
PSR-BH & Clock shift & $D_0(r_s/r)^2$ & N/A & Measurable \\
PSR-BH & Orbital precession & $D_0(v/c)^2$ & N/A & Measurable \\
\bottomrule
\end{tabular}
\end{center}

\textbf{Verdict:} Binary pulsars provide two types of DFD tests:
\begin{enumerate}
\item \textbf{Null test:} Confirm zero scalar dipole radiation (distinguishes DFD from other scalar-tensor theories). Achievable now, improved by SKA.
\item \textbf{Positive test:} Detect disformal orbital correction $\sim D_0(v/c)^4$. Requires SKA + suitable system (pulsar-BH or tighter DNS).
\end{enumerate}


\newpage
%=============================================================================
\section{Agent 5: $D_0$ via Clock Comparisons}
\label{sec:agent5}
%=============================================================================

\subsection{Mandate}

Can precision clock comparisons near massive objects detect $D_0 = 0.017$?

\subsection{New Literature (i38)}

\begin{enumerate}
\item \textbf{Zheng et al.\ (2023).} ``Lab-based test of gravitational redshift with a miniature clock network.'' Nature Comms.\ 14:4886. Measures gravitational redshift over centimeter height differences using chip-scale clocks. Fractional frequency uncertainty: $\sim 10^{-18}$. \textbf{Grade: A. State-of-the-art clock precision.}

\item \textbf{Robinson et al.\ (2025).} ``Clock Precision beyond the Standard Quantum Limit at $10^{-18}$.'' PRL. Achieves $1.1 \times 10^{-18}$ fractional frequency precision with spin-squeezed strontium lattice clock. \textbf{Grade: A. Cutting-edge.}

\item \textbf{Takamoto et al.\ (2025).} ``Verification of Gravitational Redshift Using Transportable Optical Lattice Clocks.'' 450 km baseline comparison. Uncertainty equivalent to 27 cm height. Fractional precision: $\sim 3 \times 10^{-18}$. \textbf{Grade: A. Long-distance demonstration.}

\item \textbf{Kolkowitz et al.\ (2024).} ``Artificial Precision Timing Array.'' arXiv:2401.13668. Clock satellites as decihertz GW detectors. Fractional timing: $10^{-18}$--$10^{-19}$ projected. \textbf{Grade: B. Future concept.}

\item \textbf{Ni (2024).} ``Effect of gravitational redshift on space-based GW detection with clock comparison.'' Examines systematic effects from gravitational redshift in space-based clock comparisons. \textbf{Grade: B. Systematic analysis.}
\end{enumerate}

\subsection{DFD's Clock Prediction}

\subsubsection{Disformal gravitational redshift}

In DFD, the physical clock rate is determined by the Jordan-frame metric:
\begin{equation}
\frac{d\tau_{\rm phys}}{dt} = \sqrt{-\tilde{g}_{00}} = e^{\psi}\sqrt{1 + D(\dot{\psi})^2/e^{2\psi}}\sqrt{-g_{00}}.
\end{equation}

For a static clock in a Newtonian potential $\Phi_N$:
\begin{equation}
\frac{d\tau_{\rm phys}}{dt} \approx \left(1 + \frac{\Phi_N}{c^2}\right)\left(1 + \frac{\psi}{1}\right)\left(1 + \frac{D(\nabla\psi)^2}{2}\right).
\end{equation}

In the GR regime ($\chi \gg 1$): $\psi \approx \Phi_N/c^2$ and $D(\nabla\psi)^2 = D_0\chi\Lambda_\psi^4/c^4 = D_0(\Phi_N/c^2)'^2/\Lambda_\psi^4 \times \Lambda_\psi^4 = D_0(a_{\rm grav}/c)^2/\Lambda_\psi^4 \times \Lambda_\psi^4$...

Simplifying: the disformal correction to the redshift between two clocks at heights $h_1$ and $h_2$ is:
\begin{equation}
\frac{\delta z_{\rm DFD}}{z_{\rm GR}} \sim D_0\left(\frac{\Phi_N}{c^2}\right) \sim D_0 \times 10^{-9} \quad \text{(Earth's surface)}.
\end{equation}

This gives $\delta z/z \sim 0.017 \times 10^{-9} = 1.7 \times 10^{-11}$.

The GR gravitational redshift on Earth over height $\Delta h$ is:
\begin{equation}
z_{\rm GR} = \frac{g\Delta h}{c^2} = 1.1 \times 10^{-16}\text{ per meter}.
\end{equation}

DFD correction: $\delta z \sim 1.7 \times 10^{-11} \times 1.1 \times 10^{-16}/\text{m} = 1.9 \times 10^{-27}/\text{m}$.

Current clock precision: $\sim 10^{-18}$. Required precision for DFD detection: $\sim 10^{-27}$. \textbf{Nine orders of magnitude away.}

\subsubsection{Near a neutron star}

For a neutron star ($\Phi_N/c^2 \sim 0.2$, $a_{\rm surf} \sim 10^{12}$ m/s$^2$):
\begin{equation}
\frac{\delta z_{\rm DFD}}{z_{\rm GR}} \sim D_0 \times \Phi_N/c^2 \sim 0.017 \times 0.2 \sim 3 \times 10^{-3}.
\end{equation}

A 0.3\% deviation from GR gravitational redshift at a neutron star surface. This is in principle detectable via spectral line measurements, but the systematic uncertainties from neutron star atmosphere modeling are $\sim 10\%$, overwhelming the DFD signal.

\subsubsection{Near a black hole}

At $r = 3r_s$ (ISCO): $\Phi_N/c^2 \sim 1/6$, $\psi \sim 1/6$. Disformal correction:
\begin{equation}
\delta z_{\rm DFD}/z_{\rm GR} \sim D_0/6 \sim 3 \times 10^{-3}.
\end{equation}

This could be probed by EHT photon ring measurements or X-ray spectroscopy of accretion disks. The Fe K$\alpha$ line profile depends on the spacetime geometry near the ISCO. A 0.3\% modification of the redshift at the ISCO would shift the line centroid by $\sim 20$ eV (for 6.4 keV line).

Current ATHENA/XRISM precision for Fe K$\alpha$: $\sim 2$ eV (energy resolution). \textbf{DFD's signal of 20 eV is detectable with XRISM if systematics are controlled.}

\subsection{Sensitivity Summary}

\begin{center}
\begin{tabular}{llll}
\toprule
\textbf{Method} & \textbf{$D_0$ sensitivity} & \textbf{Status} & \textbf{Timeline} \\
\midrule
Earth clocks & $D_0 < 10^{9}$ & Hopeless & --- \\
Space clocks (ISS) & $D_0 < 10^{6}$ & Hopeless & --- \\
NS redshift & $D_0 < 0.1$ & Possible & 2030s \\
BH Fe K$\alpha$ (XRISM) & $D_0 < 0.005$ & \textbf{Promising} & 2025--2030 \\
EHT photon ring & $D_0 < 0.01$ & Possible & 2030s \\
\bottomrule
\end{tabular}
\end{center}

\textbf{Verdict:} Terrestrial and space clocks cannot measure $D_0$. The best clock-based prospect is \textbf{X-ray spectroscopy of the Fe K$\alpha$ line with XRISM/ATHENA}, probing the disformal correction to the redshift at the ISCO. This is competitive with GW echoes.

\textbf{Status:} P-BH4 (disformal clock shift near BH): \YELLOW\ $\to$ \YELLOW\ (still requires breakthrough in systematics).


\newpage
%=============================================================================
\section{Agent 12: Experimental Status Update}
\label{sec:agent12}
%=============================================================================

\subsection{Mandate}

Update the status of ALL experiments relevant to quantum DFD. Any new results since i37?

\subsection{New Literature (i38)}

\begin{enumerate}
\item \textbf{AION Collaboration (2025).} ``A Prototype Atom Interferometer to Detect Dark Matter and Gravitational Waves.'' arXiv:2504.09158. AION tabletop prototype operates at the Standard Quantum Limit using $^{87}$Sr clock transition. Demonstrates noise rejection between two separated interferometers. AION-10 (10m baseline) under construction at Oxford; operation in 2--3 years. \textbf{Grade: A. Key milestone for mid-frequency GW detection.}

\item \textbf{Bose et al.\ (2026).} ``Evolution of tripartite entanglement in three-qubit QGEM.'' Sci.\ Rep.\ 16:44184. Three-qubit extension of BMV. Shows tripartite entanglement is more robust against decoherence than bipartite. Proposed as improved test of quantum gravity. \textbf{Grade: A. Directly relevant to DFD's BMV prediction.}

\item \textbf{Das, Marsh \& Bhatt (2025).} ``Using entanglement to test whether gravity is quantum just got more complicated.'' phys.org report. Shows that classical gravity can generate entanglement in some parameter regimes, complicating the BMV interpretation. If gravitons are ``heavy mediators,'' entanglement may not be detected even if gravity is quantum. Proposes dynamical fidelity susceptibility as improved probe. \textbf{Grade: A. Important complication for BMV.}

\item \textbf{IIT Bombay (2025).} Proposes dynamical fidelity susceptibility as a more sensitive probe of quantum gravity than BMV entanglement witnesses. Can distinguish quantum from classical gravity even for heavy mediators. \textbf{Grade: A. New experimental concept.}

\item \textbf{Wang et al.\ (2025).} LVK O3 echo search. No detection. $A_{\rm echo}/A_{\rm merger} < 0.02$. \textbf{Grade: A. Current bound relevant to $D_0$.}
\end{enumerate}

\subsection{Experiment-by-Experiment Status}

\subsubsection{BMV / QGEM Entanglement}

\textbf{Status:} No experiment has yet achieved the required mass ($\sim 10^{-14}$ kg) and coherence time ($\sim 1$ s) simultaneously. Current state-of-the-art: nanogram-scale oscillators with coherence times $\sim 10$ ms (Vienna, Delft groups).

\textbf{DFD relevance:} DFD predicts $\sim 2200\times$ enhancement of gravitational entanglement in the MOND regime ($a < a_0$). This means the BMV experiment could work with $\sim 50\times$ lighter masses or $\sim 50\times$ shorter coherence times than the standard quantum gravity prediction. However, the MOND regime requires sub-$a_0$ accelerations, which is challenging in a lab setting.

\textbf{Complication (new):} Das et al.\ (2025) show that classical gravity can mimic entanglement in some regimes. The DFD prediction of $2200\times$ enhancement is robust against this concern because the enhancement is a \emph{distinctive signature} that classical gravity cannot produce --- the enhancement is specifically tied to the MOND interpolation function.

\textbf{Timeline:} First BMV-type experiments: 2030--2035. MOND-enhanced BMV: requires space-based platform (2040+).

\subsubsection{AION Atom Interferometer}

\textbf{Status:} Tabletop prototype operational. AION-10 (10m) under construction at Oxford; planned operation 2028--2029. AION-100 (100m) planned for 2030s. AION-km (1km underground) long-term goal.

\textbf{DFD relevance:} AION probes the mid-frequency GW band ($\sim 0.01$--$10$ Hz). DFD predicts no modification of GW propagation (gravitons travel at $c$ exactly). AION's primary value for DFD is:
\begin{itemize}
\item Testing the equivalence principle at quantum level.
\item Potential dark matter searches (DFD has no dark matter particle, so null result supports DFD).
\item Atom-interferometric measurement of $G$ in low-acceleration environments.
\end{itemize}

\textbf{Timeline:} AION-10 first science: 2029. AION-100: 2035.

\subsubsection{LVK GW Echoes}

\textbf{Status:} O3 data analyzed. No echoes detected. O4 ongoing (to end 2025). O5 planned 2027--2029 with $\sim 2\times$ improved sensitivity.

\textbf{DFD relevance:} DFD predicts echo amplitude $A_{\rm echo}/A_{\rm merger} \sim D_0 = 0.017$. Current bound: $< 0.02$. O5 will probe $D_0 \sim 0.005$.

\textbf{Verdict:} \textbf{Most promising near-term test.} O5 will either detect DFD echoes or exclude $D_0 > 0.005$.

\subsubsection{LiteBIRD ($r$ measurement)}

\textbf{Status:} Launch planned 2032. Target sensitivity: $\sigma_r \approx 0.001$.

\textbf{DFD relevance:} i38 Agent 2 has shown that DFD's correction to $r$ is $\sim 10^{-38}$, completely undetectable. The i37 prediction of $r_{\rm DFD}/r_{\rm GR} \approx 0.85$ was a frame-confusion error. LiteBIRD is \emph{not} a useful test of DFD.

\textbf{Status:} P-INF1 ($r$ suppression): \RED\ (correction too small).

\subsubsection{SKA Binary Pulsars}

\textbf{Status:} SKA Phase 1 construction ongoing. First pulsars expected 2027--2028. Full sensitivity 2030+.

\textbf{DFD relevance:} Zero scalar dipole radiation (null test). Disformal orbital correction at $\delta P/P \sim 10^{-14}$ for DNS systems.

\textbf{Timeline:} First DFD-relevant constraints: 2030.

\subsubsection{XRISM Fe K$\alpha$ Spectroscopy}

\textbf{Status:} XRISM operational since 2023. Resolve $\sim 5$ eV (microcalorimeter). Observing Cyg X-1, GRS 1915+105.

\textbf{DFD relevance:} Disformal modification of Fe K$\alpha$ line at ISCO: $\delta E \sim 20$ eV. This is $\sim 4\times$ the instrumental resolution. \textbf{Detectable in principle.}

\textbf{Challenge:} Systematic effects from disk winds, ionization structure, and spin determination dominate at $\sim 100$ eV level. Needs careful modeling.

\textbf{Timeline:} First DFD-relevant analysis possible now with existing XRISM data. But systematics need $\sim 5$ years of calibration work.

\subsection{Master Experimental Timeline}

\begin{center}
\begin{longtable}{lllll}
\toprule
\textbf{Experiment} & \textbf{Observable} & \textbf{DFD prediction} & \textbf{Sensitivity} & \textbf{When} \\
\midrule
\endhead
LVK O5 & GW echoes & $A/A_0 \sim 0.017$ & 0.005 & 2027--29 \\
SKA Phase 1 & Dipole radiation & 0 (exact) & $10^{-5}$ & 2030 \\
SKA Phase 1 & $\delta P/P$ & $1.7\times10^{-14}$ & $10^{-14}$ & 2032 \\
XRISM & Fe K$\alpha$ shift & $\delta E \sim 20$ eV & $\sim 5$ eV & 2025--30 \\
AION-10 & EP test & Standard & Improved & 2029 \\
BMV (lab) & Entanglement & $2200\times$ enhanced & TBD & 2030--35 \\
LiteBIRD & $r$ suppression & $\sim 10^{-38}$ & $10^{-3}$ & 2032 \\
\bottomrule
\end{longtable}
\end{center}

\textbf{Priority ranking for DFD tests:}
\begin{enumerate}
\item \textbf{GW echoes (LVK O5)} --- most imminent, probes $D_0$ directly.
\item \textbf{SKA dipole radiation} --- zero prediction is sharp and distinctive.
\item \textbf{XRISM Fe K$\alpha$} --- probes near-horizon disformal geometry.
\item \textbf{BMV entanglement} --- probes quantum gravity + MOND enhancement.
\item \textbf{SKA orbital decay} --- probes disformal coupling at PN level.
\end{enumerate}

\end{document}
